%%%%%%%%%%%%%%%%%%%%%%%%%%%%%%%%%%%%%%%%%%%%%%%%%%%%%%%%%%%%%%%%%%%%%%%%%%%%%%%
%                                Basic Packages                               %
%%%%%%%%%%%%%%%%%%%%%%%%%%%%%%%%%%%%%%%%%%%%%%%%%%%%%%%%%%%%%%%%%%%%%%%%%%%%%%%

\usepackage{amsmath, amsthm, amssymb, amsfonts}
\usepackage{thmtools}
\usepackage{graphicx}
\usepackage{setspace}
\usepackage[hmargin=2.6cm]{geometry} %内部数值更改页边距
\usepackage{float}
\usepackage[colorlinks=true, citecolor=blue, linkcolor=blue, urlcolor=blue]{hyperref}
\usepackage[utf8]{inputenc}
\usepackage[english]{babel}
\usepackage{framed}
\usepackage[dvipsnames]{xcolor}
\usepackage{tcolorbox}
\usepackage{indentfirst}
\usepackage{tikz}
\usepackage{color}
\usepackage{pifont}
\usepackage{listings}
\usepackage{pdfpages}
\usepackage{enumerate}
\usepackage{mathrsfs}
\usepackage{fancyhdr}
\usepackage{lastpage}
\usepackage{layout}
\usepackage{lastpage}
\usepackage{fontawesome}
\usepackage{adjustbox}
\usepackage[ruled,longend,linesnumbered]{algorithm2e}
\usepackage{cite}
\usepackage{lipsum}
\usepackage{commutative-diagrams}
\usepackage{appendix}
\usepackage{pgfornament}
\usepackage{pgfplots}
\usepackage{titlesec}
\usepackage{booktabs}
\usetikzlibrary{shadows}

\tcbuselibrary{many}

\colorlet{LightGray}{White!90!Periwinkle}
\colorlet{LightOrange}{Orange!15}
\colorlet{LightGreen}{Green!15}
\colorlet{White}{White!}

\definecolor{mygreen}{RGB}{56, 140, 70}
\definecolor{mypurple}{RGB}{197, 92, 212}
\definecolor{nordAurora1}{HTML}{BF616A}
\definecolor{lemma}{HTML}{983b0f}
\definecolor{prop}{HTML}{191971}
\definecolor{definition}{HTML}{88D6D1}
\definecolor{myblue}{RGB}{0,127,255}
\definecolor{blue2}{HTML}{99ebff}
\definecolor{lightergray}{gray}{0.99}
\definecolor{main}{RGB}{0,166,82}
\definecolor{third}{RGB}{0,174,247}
\definecolor{partcolor}{HTML}{3352cc}

\colorlet{example}{mygreen!85!black}
\colorlet{claim}{mypurple!90!black}

\numberwithin{equation}{section}
\numberwithin{figure}{section}
\newcommand{\HRule}[1]{\rule{\linewidth}{#1}}
\newcommand{\solutionname}{Solution}

\declaretheoremstyle[name=\colorbox{Orange!80}{\textcolor{White}{Theorem}},
headfont=\bfseries\color{Orange!90},
bodyfont=\itshape,]{thm}
\declaretheorem[style=thm,numberwithin=section]{theorem}
\tcolorboxenvironment{theorem}{breakable,
colback=LightOrange!85,
colframe=Orange!80}

\declaretheoremstyle[name=\textcolor{lemma}{Lemma},
headfont=\bfseries\color{lemma},
bodyfont=\itshape,]{lem}
\declaretheorem[style=lem,numberwithin=section]{lemma}
\tcolorboxenvironment{lemma}{breakable,
colback=LightOrange!37,
colframe=lemma}

\declaretheoremstyle[name=\colorbox{myblue}{\textcolor{White}{Definition}},
headfont=\bfseries\color{myblue},
bodyfont=\itshape,]{def}
\declaretheorem[style=def,numberwithin=section]{definition}
\tcolorboxenvironment{definition}{breakable,
colback=blue2!50,
colframe=myblue}

\declaretheoremstyle[name=\colorbox{mygreen!100!black}{\textcolor{White}{Corollary}},
headfont=\bfseries\color{mygreen!100!black},
bodyfont=\itshape,]{coro}
\declaretheorem[style=coro,numberwithin=section]{corollary}
\tcolorboxenvironment{corollary}{breakable,
colback=LightGreen!65,
colframe=mygreen!100!black}

\declaretheoremstyle[name=\textcolor{prop}{Proposition},
headfont=\bfseries\color{prop},
bodyfont=\itshape,]{pro}
\declaretheorem[style=pro,numberwithin=section]{proposition}
\tcolorboxenvironment{proposition}{breakable,
colback=LightGray,
colframe=prop}

\declaretheoremstyle[name=\colorbox{claim}{\textcolor{White}{Claim}},
headfont=\bfseries\color{claim},
bodyfont=\itshape,]{claim}
\declaretheorem[style=claim,numberwithin=section]{claim}
\tcolorboxenvironment{claim}{breakable,
colback=nordAurora1!11!White,
colframe=claim}

\declaretheoremstyle[name=\textcolor{red}{Problem},
headfont=\bfseries\color{red},
bodyfont=\upshape,]{problem}
\declaretheorem[style=problem,numberwithin=section]{problem}
\tcolorboxenvironment{problem}{breakable,
colback=White}

\declaretheoremstyle[name=\textcolor{red}{Question},
headfont=\bfseries\color{red},
bodyfont=\upshape,]{que}
\declaretheorem[style=que,numberwithin=section]{question}
\tcolorboxenvironment{question}{breakable,
colback=White}


\declaretheoremstyle[name=\textcolor{example}{Example},
headfont=\bfseries\color{example},
bodyfont=\upshape,]{exm}
\declaretheorem[style=exm,numberwithin=section]{example}
\tcolorboxenvironment{example}{breakable,
colback=LightGreen!65,
colframe=example}

\let\proof\relax
\let\endproof\relax
\newenvironment{proof}{
  \par
  \vspace{.5\baselineskip}
  \noindent\textbf{\color{Orange!90!black} Proof.\;}
  \color{black!88}
}{
  \hfill$\blacksquare$\quad
  \par\vspace{.5\baselineskip}
}

\newenvironment{solution}{\par\vspace{.5\baselineskip}\noindent\textbf{\color{main}\solutionname:} \itshape}{\color{main}\hfill\ding{118}\par\vspace{.5\baselineskip}}

\newenvironment{remark}{\par\vspace{.5\baselineskip}\noindent\textbf{\color{Orange!90!black}Remark.\;}}{\par\vspace{.5\baselineskip}}

\newenvironment{lanse}{
  \begin{tcolorbox}[enhanced, colback=LightOrange!30, boxrule=0pt, frame hidden, breakable,
  borderline west={0.7mm}{0.1mm}{lemma}]
}{\end{tcolorbox}}
\newtheorem*{notationev}{\color{lemma} Notation}
\newenvironment{notation}{\begin{lanse}\begin{notationev}}{\hfill\ding{45}\end{notationev}\end{lanse}}

\newenvironment{exer}{
  \begin{tcolorbox}[enhanced, colback=LightGreen!65, boxrule=0pt, frame hidden, breakable,
  borderline west={0.7mm}{0.1mm}{example}]
}{\end{tcolorbox}}
\newtheorem{exerciseev}{\color{example} Exercise}[section]
\newenvironment{exercise}{\begin{exer}\begin{exerciseev}}{\end{exerciseev}\end{exer}}

\newcounter{property}[section]
\renewcommand{\theproperty}{\thesection.\arabic{property}}

\newenvironment{property}{\refstepcounter{property}\par\vspace{.5\baselineskip}\noindent\textbf{\color{third} Property \theproperty} \itshape}{\color{third}\hfill\ding{168}\par\vspace{.5\baselineskip}}

% 猜想环境
\newcounter{conj}[section]
\setcounter{conj}{0}
\renewcommand{\theconj}{\thesection.\arabic{conj}}

\newtcolorbox{conjbox}[3][]{%
    arc=0mm,
    breakable,drop fuzzy shadow=black!20!white,
    enhanced,
    colback=lightergray,%main!3
    boxrule=0pt,
    top=8mm, %分隔垂直-开始文本
    % enlarge top by=\baselineskip/2+1mm,
    % enlarge top at break by=0mm,pad at break=2mm,
    fontupper=\normalsize,
    step and label={conj}{#3},%标签
    overlay unbroken = {%
        %垂直线段%更改有效
        %\draw[color=gray!5,line width=3pt] ([xshift=2pt] frame.north west)--([xshift=2pt] frame.south west);
        %底部装饰矩形%更改有效
        \fill[fill=red!80!black] (.9\textwidth,0.2cm) rectangle (\textwidth,-0.0cm);
        % 图像框示例
        %\node[rectangle,
        %         text=black, ,fill=ALamarillo!20,
        %         inner sep=1mm,anchor=west,font=\small\bf\sffamily] at ([xshift=-10.3pt,yshift=-3.2mm]frame.north west)%
        %{ };
        %框编号和描述%更改有效
        \node[rectangle, %opacity=.3,
        text=black, drop shadow={opacity=.3, shadow xshift=0.1cm},
        inner sep=1.5mm,fill=red!80!black,
        anchor=west,
        font=\normalsize] at ([xshift=0cm,yshift=-3.mm]frame.north west)%
        {\bfseries  \color{white}Conjecture \theconj~#2};%标题
    },
    % fin overlay
    overlay first  = {%
        %标题框编号和描述%跨页时生效
        \node[rectangle, %opacity=.3,
        text=black, drop shadow={opacity=.3, shadow xshift=0.1cm},
        inner sep=1.5mm,fill=Emerald,
        anchor=west,
        font=\normalsize] at ([xshift=0cm,yshift=-3.mm]frame.north west)%
        {\bf \color{white} Conjecture \theconj.~#2};
    },
    % overlay
    %边缘改变%跨页时生效
    overlay middle={\draw[color=lightergray,line width=3pt] ([xshift=3pt] frame.north west)-- ([xshift=2pt] frame.south west);
    },
    overlay last={\draw[color=lightergray,line width=3pt] ([xshift=3pt] frame.north west)-- ([xshift=2pt] frame.south west);
        %底部装饰矩形%跨页时生效
        \fill[fill=red!80!black] (.9\textwidth,0.2cm) rectangle (\textwidth,-0.0cm); }
    #1
}

\NewDocumentEnvironment{conjecture}{O{} O{} O{}}{\smallskip\begin{conjbox}{#1}{#2}%
        #3}{\end{conjbox} }

%悖论环境
\newcounter{Para}[section]
\setcounter{Para}{0}
\renewcommand{\thePara}{\thesection.\arabic{Para}}

\newtcolorbox{parabox}[3][]{%
    arc=0mm,
    breakable,drop fuzzy shadow=black!20!white,
    enhanced,
    colback=lightergray,%main!3
    boxrule=0pt,
    top=8mm, %分隔垂直-开始文本
    % enlarge top by=\baselineskip/2+1mm,
    % enlarge top at break by=0mm,pad at break=2mm,
    fontupper=\normalsize,
    step and label={Para}{#3},%标签
    overlay unbroken = {%
        %垂直线段%更改有效
        %\draw[color=gray!5,line width=3pt] ([xshift=2pt] frame.north west)--([xshift=2pt] frame.south west);
        %底部装饰矩形%更改有效
        \fill[fill=red!80!black] (.9\textwidth,0.2cm) rectangle (\textwidth,-0.0cm);
        % 图像框示例
        %\node[rectangle,
        %         text=black, ,fill=ALamarillo!20,
        %         inner sep=1mm,anchor=west,font=\small\bf\sffamily] at ([xshift=-10.3pt,yshift=-3.2mm]frame.north west)%
        %{ };
        %框编号和描述%更改有效
        \node[rectangle, %opacity=.3,
        text=black, drop shadow={opacity=.3, shadow xshift=0.1cm},
        inner sep=1.5mm,fill=red!80!black,
        anchor=west,
        font=\normalsize] at ([xshift=0cm,yshift=-3.mm]frame.north west)%
        {\bfseries  \color{white}Paradox \thePara~#2};%标题
    },
    % fin overlay
    overlay first  = {%
        %标题框编号和描述%跨页时生效
        \node[rectangle, %opacity=.3,
        text=black, drop shadow={opacity=.3, shadow xshift=0.1cm},
        inner sep=1.5mm,fill=Emerald,
        anchor=west,
        font=\normalsize] at ([xshift=0cm,yshift=-3.mm]frame.north west)%
        {\bf \color{white} Paradox \thePara.~#2};
    },
    % overlay
    %边缘改变%跨页时生效
    overlay middle={\draw[color=lightergray,line width=3pt] ([xshift=3pt] frame.north west)-- ([xshift=2pt] frame.south west);
    },
    overlay last={\draw[color=lightergray,line width=3pt] ([xshift=3pt] frame.north west)-- ([xshift=2pt] frame.south west);
        %底部装饰矩形%跨页时生效
        \fill[fill=red!80!black] (.9\textwidth,0.2cm) rectangle (\textwidth,-0.0cm); }
    #1
}

\NewDocumentEnvironment{paradox}{O{} O{} O{}}{\smallskip\begin{parabox}{#1}{#2}%
        #3}{\end{parabox} }

\setstretch{1.2}
\geometry{
    hmargin= 2.5cm,
    textheight=9in,
    top=1in,
    headheight=14pt,
    headsep=25pt,
    footskip=30pt
}

%%%%%%%%%%%%%%%%%%%%%%%%%%%%%%%%%%%%%%%%%%%%%%%%%%%%%%%%%%%%%%%%%%%%%%%%%%%%%%%
%                                 Environments                                %
%%%%%%%%%%%%%%%%%%%%%%%%%%%%%%%%%%%%%%%%%%%%%%%%%%%%%%%%%%%%%%%%%%%%%%%%%%%%%%%

\DeclareStringOption[Blue]{color}
\DeclareVoidOption{green}{\ekv{color=green}}
\DeclareVoidOption{cyan}{\ekv{color=cyan}}
\DeclareVoidOption{Blue}{\ekv{color=Blue}}
\DeclareVoidOption{sakura}{\ekv{color=sakura}}
\DeclareVoidOption{black}{\ekv{color=black}}
\DeclareVoidOption{brown}{\ekv{color=brown}}

%%%%%%%%%%%%%%%%%%%%%%%%%%%%%%%%%%%%%%%%%%%%%%%%%%%%%%%%%%%%%%%%%%%%%%%%%%%%%%%
%                                 Code blocks                                 %
%%%%%%%%%%%%%%%%%%%%%%%%%%%%%%%%%%%%%%%%%%%%%%%%%%%%%%%%%%%%%%%%%%%%%%%%%%%%%%%

\lstset{backgroundcolor = \color{LightOrange!20},
rulesepcolor=\color{red!20!green!20!blue!20},
showspaces=false,
showstringspaces=false,}
\lstdefinestyle{lfonts}{
  basicstyle   = \small\ttfamily,
  stringstyle  = \color{green!70!black},
  keywordstyle = \color{blue!68!white}\bfseries,
  commentstyle = \color{olive},
}
\lstdefinestyle{lnumbers}{
  numbers     = left,
  numberstyle = \tiny,
  numbersep   = 2em,
  firstnumber = 1,
  stepnumber  = 1,
}
\lstdefinestyle{llayout}{
  breaklines       = true,
  tabsize          = 4,
  columns          = flexible,
}
\lstdefinestyle{lgeometry}{
  xleftmargin      = 20pt,
  xrightmargin     = 0pt,
  frame            = shadowbox,
  framesep         = \fboxsep,
  framexleftmargin = 20pt,
}
\lstdefinestyle{lgeneral}{
  style = lfonts,
  style = lnumbers,
  style = llayout,
  style = lgeometry,
}
\lstdefinestyle{python}{
    language = {Python},
    style    = lgeneral,
}